\documentclass[11pt,uplatex,dvipdfmx,a4paper]{jsarticle}

\usepackage[dvipdfmx]{graphicx}
\usepackage[dvipdfmx,colorlinks=true,linkcolor=blue,citecolor=blue]{hyperref}

\usepackage{amsmath}
\usepackage{physics2}
\usephysicsmodule{ab}

\usepackage{authblk}
\renewcommand\Authand{,\ } % 著者の区切りをカンマに変更

\title{\textbf{シンプル高速\LaTeX{}テンプレート}}
% 著者と所属は必要に応じてここで指定します。
% 単一著者の場合
% \author{山田 太郎}
% \affil{東京大学}

\author[A]{山田 太郎}
\author[AB]{鈴木 花子}
\affil[A]{東京大学}
\affil[B]{京都大学}
\date{\today}

% 参考文献スタイルはここで切り替えます。
% `unsrt` なら [1] のような番号形式、
% `apsrev4-2` なら APS 風の形式になります。
\bibliographystyle{apsrev4-2}

\begin{document}
\maketitle
\begin{abstract}
これは日本語で論文を書くためのシンプルで高速な\LaTeX{}テンプレートです。キャッシュされたプリアンブルを使用して、コンパイルを高速化しています。レイアウトは可能な限りシンプルで、参考文献スタイルは用途に応じて切り替えられます。
\end{abstract}

\section{導入}

このテンプレートは、シンプルで高速な\LaTeX{}環境を提供することを目的としています。
キャッシュされたプリアンブルを使用することで、コンパイル時間を大幅に短縮しています。
日本語の文書もそのまま書くことができます。
参考文献の引用例として、論文は~\cite{einstein1905}、
書籍は~\cite{lamport1994} のように書けます。

数式は以下のように書くことが可能です。
\begin{align}
G &= \ab(V, E).
\end{align}

\bibliography{main}

\end{document}
